\chapter{系统总体设计}
\label{cha:system_design}

\section{系统设计目标}
\label{sec:design_goals}

SkillRank的系统设计目标主要包括以下几个方面：

\begin{itemize}
    \item 支持大规模Skill数据的统一接入与管理；
    \item 支持多维特征建模与综合价值评估；
    \item 支持面向用户查询场景的筛选、排序与详情展示；
    \item 保证系统架构清晰，便于后续扩展与维护。
\end{itemize}

\section{系统总体架构}
\label{sec:system_architecture}

SkillRank采用分层式云端系统架构。数据接入层负责从外部平台或数据源收集Skill相关信息，并完成原始数据的拉取与同步；数据处理层负责元数据清洗、字段统一和特征抽取；价值评估层根据预设指标体系计算Skill综合评分；服务接口层向前端提供查询、筛选、排序和详情访问能力；前端展示层面向最终用户提供可视化交互界面。各模块分工清晰，能够较好支撑Skill数据从采集到展示的完整闭环。

系统架构可分为以下五层：

\begin{enumerate}
    \item \textbf{数据接入层}：负责从GitHub、社区平台等数据源获取Skill原始信息
    \item \textbf{数据处理层}：完成数据清洗、标准化和特征提取
    \item \textbf{价值评估层}：基于多维特征计算Skill综合评分
    \item \textbf{服务接口层}：提供RESTful API支持查询、筛选和排序
    \item \textbf{前端展示层}：提供用户交互界面和可视化展示
\end{enumerate}

\section{系统功能模块设计}
\label{sec:functional_modules}

\subsection{Skill数据接入模块}
\label{subsec:data_ingestion}

数据接入模块负责从多个数据源获取Skill信息。主要数据来源包括GitHub仓库、技能市场平台和社区分享。数据拉取采用定期同步机制，每日更新一次基础信息，每周更新一次统计指标。

数据清洗规则包括：去除重复Skill、标准化字段格式、处理缺失值、过滤无效数据。异常处理方式采用日志记录和人工审核相结合的方式，确保数据质量。

\subsection{Skill元数据管理模块}
\label{subsec:metadata_management}

元数据管理模块存储Skill的核心信息，包括：Skill名称、描述、作者、来源链接、创建时间、更新时间、分类标签、依赖关系等。

对于标签、作者、描述等信息，系统进行标准化处理：统一标签命名规范、去重作者信息、规范化描述格式。通过哈希值和内容相似度进行去重判断，确保数据一致性。

\subsection{价值评估与排序模块}
\label{subsec:value_evaluation}

价值评估与排序模块是SkillRank系统的核心。本模块的设计基于对Agent Skill特性的深入分析，将Skill价值评估问题分解为两个相互独立但互补的维度：需求匹配维度与内容质量维度。

\subsubsection{问题建模}

在Agent Skill生态中，每个Skill可以形式化表示为一个二元组：

\begin{equation}
    \label{eq:skill_structure}
    Skill = \langle Description, Content \rangle
\end{equation}

其中，Description作为触发器（Trigger），决定Skill是否被匹配到特定需求；Content作为信息载体，决定Skill被使用后对Agent输出质量的实际影响。这种二元结构类似于传统信息检索中的"索引-文档"关系，但在Agent场景下具有独特的语义特征。

借鉴PageRank算法的核心思想，本文提出基于内容亲缘关系的Skill价值评估机制。PageRank通过网页间的超链接关系评估网页价值，其核心假设是：被更多网页链接的网页更重要，被高质量网页链接的网页更重要。类比到Skill场景，本文提出如下假设：

\begin{itemize}
    \item 假设1：Content相似度高的Skill之间存在"亲缘关系"，反映了某种被广泛认可的解法模式
    \item 假设2：被更多Skill的Content引用或模仿的Skill，代表了更基础、更通用的能力封装
    \item 假设3：与高质量Skill具有高Content相似度的Skill，其自身质量也更可能较高
\end{itemize}

基于上述假设，本文将Skill价值评估问题建模为需求空间中的"地形分布"问题。在以Description embedding为坐标的需求空间中，每个Skill占据特定位置，其Content质量决定该位置的"高度"。高质量Skill在需求空间中形成"高峰"，而相似需求区域内的Skill质量分布构成连续的"地形"。

\subsubsection{评分维度设计}

基于上述建模，本系统的评分机制包含以下维度：

\begin{enumerate}
    \item \textbf{热度特征（Popularity）}：包括GitHub Star数、Fork数、下载量、浏览量等指标，反映Skill的社区认可度
    \item \textbf{活跃度特征（Activity）}：包括最近更新时间、提交频率、维护状态等，反映Skill的生命力
    \item \textbf{内容质量特征（Quality）}：包括描述完整性、文档丰富度、示例数量、代码质量等，反映Skill的可用性
    \item \textbf{可复用性特征（Reusability）}：包括标签规范性、参数化程度、适用场景广度、依赖复杂度等，反映Skill的通用性
    \item \textbf{生态反馈特征（Feedback）}：包括被引用次数、收藏数、社区讨论热度等，反映Skill的实际使用情况
    \item \textbf{内容亲缘特征（Content Affinity）}：基于Content相似度构建的亲缘图，通过图算法计算Skill在生态中的结构重要性
\end{enumerate}

综合评分公式采用加权求和方式：

\begin{equation}
    \label{eq:skill_score}
    Score(skill) = \sum_{i=1}^{6} w_i \cdot F_i(skill)
\end{equation}

其中，$F_i$表示第$i$个维度的归一化特征值，$w_i$为对应权重，满足$\sum_{i=1}^{6} w_i = 1$。

\subsubsection{内容亲缘图构建}

内容亲缘特征的计算基于Skill间的Content相似度图。对于任意两个Skill $s_i$和$s_j$，其Content相似度定义为：

\begin{equation}
    \label{eq:content_similarity}
    sim(s_i, s_j) = \frac{emb(Content_i) \cdot emb(Content_j)}{\|emb(Content_i)\| \cdot \|emb(Content_j)\|}
\end{equation}

其中，$emb(\cdot)$表示Content的向量表示（通过预训练语言模型获得）。基于相似度矩阵，构建无向加权图$G = (V, E)$，其中$V$为Skill集合，边权重为Content相似度。在此图上应用图中心性算法（如Degree Centrality、Betweenness Centrality或改进的PageRank），计算每个Skill的结构重要性得分。

排序结果根据综合分数从高到低生成，支持多条件筛选后的二次排序，确保结果的可解释性和合理性。

\subsection{查询与展示模块}
\label{subsec:query_display}

查询与展示模块为用户提供灵活的检索和浏览功能，主要包括：

\begin{itemize}
    \item \textbf{关键词搜索}：支持对Skill名称、描述、标签进行全文检索
    \item \textbf{标签筛选}：支持按分类标签、技术栈、应用场景等维度过滤
    \item \textbf{多维排序}：支持按综合分数、热度、活跃度、更新时间等排序
    \item \textbf{详情查看}：展示Skill的完整信息、关键指标和价值评分构成
\end{itemize}

\section{数据库设计}
\label{sec:database_design}

系统采用关系型数据库存储Skill相关数据，核心表结构如下：

\textbf{skills表}：保存Skill的基础信息，包括Skill名称、描述、来源链接、作者、创建时间、更新时间、分类标签和状态等字段。

\textbf{skill\_tags表}：存储Skill与标签的多对多关系，支持灵活的标签管理和检索。

\textbf{skill\_metrics表}：记录Skill的各项统计指标，如Star数、Fork数、下载量、浏览量等，支持历史数据追踪。

\textbf{skill\_scores表}：存储Skill的价值评分结果，包括各维度得分和综合得分，便于快速排序查询。

\textbf{authors表}：管理Skill作者信息，包括作者名称、联系方式、贡献统计等。

\section{接口设计}
\label{sec:api_design}

系统提供RESTful API接口，支持前端和第三方应用访问。核心接口如下：

\begin{table}[!htbp]
    \centering
    \caption{系统核心接口}
    \label{tab:api_list}
    \begin{tabular}{|l|l|l|}
        \hline
        \textbf{接口名称} & \textbf{请求方式} & \textbf{功能说明} \\ \hline
        /api/skills & GET & 获取Skill列表 \\ \hline
        /api/skills/\{id\} & GET & 获取Skill详情 \\ \hline
        /api/skills/search & GET & 按条件查询Skill \\ \hline
        /api/scores/refresh & POST & 刷新评分结果 \\ \hline
        /api/tags & GET & 获取标签列表 \\ \hline
    \end{tabular}
\end{table}

\section{本章小结}
\label{sec:design_summary}

本章从系统目标、总体架构、功能模块、数据库设计和接口设计等方面给出了SkillRank的总体设计方案。通过分层架构和模块化设计，系统能够支持大规模Skill数据的接入、组织、评估与展示，为后续系统实现提供了结构基础。
